\section{Метризуемость слабых топологий.}\label{seq:metrizability}
\begin{theorem}
    Пусть $X$ -- нормированное пространство.
    \begin{enumerate}
        \item Если $X^*$ -- сепарабельно, то слабая топология, индуцированная на шар $B_R(0)$ метризуема.
        \item Если $X$ -- сепарабельно, то *-слабая топология, индуцированная на шар $B_R^*(0)$ метризума.
    \end{enumerate}
\end{theorem}
\begin{proof}
    Случай $X = 0$ тривиален.
    Пусть $X \neq 0$.
    Тогда и $X^* \neq 0$ и, в силу сепарабельности, существует $\{f_n\} \subset S_1^*(0)$ -- всюду плотное счётное множество векторов на единичной сфере сопряжённого пространства.
    Определим $\rho\colon X \times X \to \R$ как
    \[
        \rho\colon (x, y) \mapsto \sum_{n = 1}^{\infty} 2^{-n}|f_n(x - y)|.
    \]
    Заметим, что $\rho$ корректно определена.
    То есть если $x, y \in X$, то
    \[
        \rho(x, y) \leq \sum_{n = 1}^{\infty} 2^{-n} \|f_n\|\|x - y\| \leq \|x - y\| < +\infty.
    \]
    Очевидно, что $\rho(x, y) = 0 \Longleftrightarrow \forall n \in \N f_n(x) = f_n(y)$.
    Так как $\forall g \in X^*\colon \|g\| = 1 \ \forall \epsilon > 0 \exists f_n\colon \|g - f_n\| \leq \epsilon$.
    Но, тогда
    \[
        |g(x - y)| = |g(x - y) - f(x - y)| \leq \epsilon \|x - y\|.
    \]
    Поскольку это верно для произвольного $\epsilon$, то $g(x) = g(y)$.
    А значит, по следствию из теоремы Хана-Банаха, $x = y$.
    Поскольку симметричность и неравенства треугольника очевидны, то $\rho$ -- действительно метрика.

    Рассмотрим $B_R(0)$ с $\tau_{\rho}|_{B_R(0)}$ и $\tau_{w}|_{B_R(0)}$ -- метрическая и слабая топологии соответственно, индуцированные на шар.
    Покажем то, что $\tau_w|_{B_R(0)} \subset \tau_\rho|_{B_R(0)}$.

    Пусть $x \in B_R(0), \ \epsilon > 0$ и $f \in X^*$.
    Достаточно показать, что $\exists r > 0$ такое, что
    \[
        V(x, f, \epsilon) \cap B_R(0) \supset O_r^{\rho}(x) \cap B_R(0).
    \]
    Тогда у нас для всякого $G \in \tau_w|_{B_R(0)}$ и всякой $x \in G$ существует $\{V(x, g_j, \epsilon)\}_{j = 1}^{N}$
    такие, что
    \[
        G \supset \bigcap_{j = 1}^{N} V(x, g_j, \epsilon) \supset \bigcap_{j = 1}^{N} O_{r_j}^{\rho}(x) \ni x,
    \]
    а значит каждая точка $G$ входит в него вместе с некоторой метрической окрестностью, и тогда $G \in \tau_\rho|_{B_R(0)}$.

    Возвращаясь к поиску такого $r$.
    Если $f = 0$, то $V(x, 0, \epsilon) = X$ и $r = 1$, например, подойдёт (как, собственно, и любой другой $r > 0$).
    Пусть $f \neq 0$.
    Тогда $\forall \delta > 0 \ \exists f_N\colon \left\|\frac{f}{\|f\|} - f_N\right\| < \delta$.
    Пусть $y \in O_r^{\rho}(x)$.
    Это означает, что
    \[
        \rho(x, y) < r \Longleftrightarrow \sum_{n = 1}^{\infty} 2^{-n}|f_n(x - y)| < r.
    \]
    В частности, можно заметить, что
    \[
        |f_N(x - y)| < 2^n r.
    \]
    Но
    \[
        |f(x - y)| = \|f\|\left|\frac{f}{\|f\|}(x - y) - f_N(x - y) + f_N(x - y)\right| \leq \|f\| \delta \|x - y\| + \|f\|2^n r \leq 2R\|f\|\delta 2^n\|f\|r.
    \]
    Возьмём $r = \frac{\epsilon}{3\cdot 2^n\|f\|}$ и $\delta = \frac{\epsilon}{6\|f\|R}$ и получим искомую оценку.

    Покажем обратное вложение, то есть то, что $\tau_\rho|_{B_R(0)} \subset \tau_w|_{B_R(0)}$.
    Рассмотрим $O_r^{\rho}(x)$, где $x \in B_R(0)$.
    Аналогично, достаточно показать, что существуют $N \in \N$ и $\epsilon > 0$ такие, что
    \[
        O_r^{\rho}(x) \supset \bigcap_{k = 1}^{N} V(x, f_k, \epsilon).
    \]
    Рассмотрим $y \in \bigcap_{k = 1}^{N} V(x, f_k, \epsilon)$.
    Тогда
    \[
        \rho(x, y) = \sum_{k = 1}^{\infty} 2^{-k}|f_k(x - y)| = \sum_{k = 1}^{N}2^{-k}|f_k(x - y)| + \sum_{k = N + 1}^{\infty} 2^{-k}|f_k(x - y)| \leq \epsilon + 2^{-N}\|x - y\| \leq \epsilon + 2^{1-N}R.
    \]
    Можно взять $\epsilon < \frac{r}{2}$ и $N \in \N$ такое, чтобы $2^{1-N}R < \frac{r}{2}$ и прийти к победе.


    Перейдём теперь к *-слабому случаю.
    Пусть $\{x_n\}$ -- счётное всюду плотное на единичной сфере $X$ множество.
    Определим метрику на $X^*$ как
    \[
        \rho_*(f, g) = \sum_{n = 1}^{\infty} 2^{-n}|(f - g)(x_n)| \leq \|f - g\|.
    \]
    Так как $\rho(f, g) = 0 \Longleftrightarrow f(x_n) = g(x_n)$, то $f = g$ на единичной сфере (а значит и на всём $X$).
    Симметричность и неравенство треугольника проверяются абсолютно аналогично.
    Собственно, доказательство вложения туда-сюда получаются заменой точек на функционалы и функционалов на точки.
    И дописыванию <<*>> в нужных местах.
    % TODO: дописать * в нужных местах
\end{proof}
\begin{notet}
    На самом деле, это два частных случая одного утверждения
    \begin{proposition}
        Если $X$ -- нормированное пространство, $\FF \subset X^*$ -- сепарабельное линейное пространство, то $\FF$-слабая топология на $X$ метризуема на всяком ограниченном шаре.
    \end{proposition}
    Более того, доказательство этого факта проводится абсолютно аналогично доказательству теоремы.
\end{notet}
\begin{note}
    Определённая в доказательстве теоремы метрика для $X$ не порождает слабую сходимость (если $X$ -- бесконечномерно).
    Так, для всякого $m \in \N$ верно, что коразмерность $\bigcap_{n = 1}^{m} \ker f_n$ конечна и не более $m$, то есть
    \[
        X = \bigcap_{n = 1}^{m} \ker f_n \oplus \Lin\{z_1, \ldots, z_m\},
    \]
    для некоторых $z_1, \ldots, z_m \in X$.
    В силу бесконечномерности пересечения ядер существует $x_m$ такой, что $\|x_m\| = m$.
    В то же время $\rho(x_m, 0) \leq \sum_{n = m + 1}^{\infty} 2^{-n}m \rightarrow 0, m \rightarrow \infty$.
    То есть $\{x_m\}$ $\rho$-сходится к $0$, но при этом не является слабо сходящейся: всякая слабо сходящаяся последовательность ограничена по норме, а $\{x_m\}$ -- нет.
\end{note}
\begin{proposition}\label{prop:conjugate_separability}
    Пусть $X$ -- нормированное пространство и $X^*$ -- сепарабельно.
    Тогда и $X$ является сепарабельным.
\end{proposition}
\begin{proof}
    Пусть $\{f_n\}_{n = 1}^{\infty}$ -- счётное всюду плотное множество из $X^*$.

    Так как
    \[
        \|f_n\| = \sup\limits_{\|x \| = 1} |f_n(x)|,
    \]
    то для любого $n \in \N$ существует $x_n$ единичной нормы такое, что
    \[
        \|f_n\| \leq 2|f_n(x_n)|,
    \]
    в силу определения супремума.

    Рассмотрим
    \[
        M = \left\{\sum_{k = 1}^{n} \alpha_k x_k \mid n \in \N, \alpha_k \in \Cm, \Re \alpha_k \in \Q, \Im \alpha_k \in \Q\right\}.
    \]
    Очевидно, что оно счётно (при фиксированном размере линейной комбинации $m$ множество можно отождествить с $\Q^{3m}$, так как $|\{x_n\}| \cong |\Q|$, а каждый коэффициент определяется двумя рациональными числами; счётное объединение счётных множеств -- счётно).

    Пусть $L = [M]_{\|\cdot\|}$.
    Заметим, что $L$ -- некоторое подпространство $X$.
    И действительно: $M$ замкнуто относительно сложения, а значит и $L$ замкнуто относительно сложения.
    В тоже время, так как $\Q + i\Q$ плотно в $\Cm$, то любое $\alpha \in \Cm$ можно приблизить последовательностью комплексно-рациональных чисел.
    Тогда, если $x \in L$ и $\alpha x \in L$, то $\forall \epsilon > 0$ существует $\alpha_0$ такое, что $|\alpha - \alpha_0| < \frac{\epsilon}{\|x\|}$ и существует $\{y_j\}_{j = 1}^{N} \subset \{x_n\}$ такие, что при некоторых $\{\alpha_j\}_{j = 1}^{N} \subset \Q + i\Q$ верно, что
    \[
        \left\|x - \sum_{j = 1}^{N} \alpha_j y_j\right\| \leq \frac{\epsilon}{|\alpha| + |\alpha_0| + 1}.
    \]

    А значит
    \[
        \left\|\alpha x - \sum_{j = 1}^{N}\alpha_0 \alpha_j y_j\right\| \leq |\alpha - \alpha_0|\|x\| + |\alpha_0|\left\|x - \sum_{j = 1}^{N} \alpha_j y_j \right\| \leq \epsilon + \frac{|\alpha_0|\epsilon}{|\alpha| + |\alpha_0| + 1} < 2\epsilon.
    \]
    Поскольку это верно для произвольного $\epsilon$, то $\alpha x \in L$ и $L$ -- подпространство.

    Пусть $L \neq X$.
    Тогда существует $z \in X \setminus L$.
    $L$ -- замкнутое подпространство в $X$, и по следствию из теоремы Хана-Банаха существует такой непрерывный функционал $f$, что $f|_{L} = 0$ и $f(z) = 1$.
    С другой стороны, в силу всюду плотности $\{f_n\}$ у нас $\forall \epsilon > 0$ существует $f_n$ такой, что $\|f - f_n\| < \epsilon$
    Таким образом
    \[
        \|f\| \leq \|f - f_n\| + \|f_n\| \leq \epsilon + 2|f_n(x_n)| = \epsilon + 2|(f_n - f)(x_n)| \leq \epsilon + 2\|f - f_n\|\|x_n\| \leq 3\epsilon.
    \]
    Это верно для произвольного $\epsilon$, а значит $\|f\| = 0$ и $f = 0$ -- противоречие, $f(z) \neq 0$.
    Следовательно, $X$ -- сепарабельно.
\end{proof}
\begin{note}
    Обратное неверно: $\ell^1$ -- сепарабельно, но $(\ell^1)^* \cong \ell^\infty$ -- несепарабельно.
\end{note}
